\section{논의와 한계}\label{sec:discussion}
\subsection{현재 결과가 지지하는 범위}
이미 학습된 추정기의 코너를 수정한다는 개념과 치수·대칭 활용 자체는 선행연구에 존재한다. 본 결과가 보여주는 것은 내부 영상 특징과 알려진 치수로 제한된 이동을 학습했을 때, 세 기반에서 중심부 오차가 어떻게 줄고 어떤 실패가 남는지이다. 여러 기반의 결과는 각 기반에 별도로 학습한 보정기의 적용성으로 해석해야 한다. 초기 디코더와 검출 조건까지 완전히 같게 만든 백본 치환 실험이나 하나의 가중치를 임의 추정기에 전이한 결과는 아니다.

Base 대비 이득에는 추가 보정 학습 전체의 효과가 포함된다. N0--N3 절제는 치수 분기와 회전 대응의 역할을 더 좁히지만 치수 내용과 파라미터 효과를 완전히 분리하지 못한다. 치수 입력이 후보 점수에 영향을 준다는 것과 최종 코너가 물리 형상을 일관되게 복원한다는 것도 다르다. 더 강한 형상 일관성이나 초기 추정기 직접 치수 입력의 이득은 이 후단 보정 결과로 대체할 수 없다.

\subsection{대안 선택과 실패의 의미}
N3에 포함된 회전 대응 감독의 추가 이득은 작거나 혼합되어 있다. N2를 더 단순한 후보로 논의할 수는 있으나, N3의 결과를 N2의 결과로 바꾸어 부를 수는 없다. PoseFix 기반 보정도 여러 중앙값에서 N3보다 우수하므로, P90 하나만으로 제안 방법을 전반적으로 우월하다고 결론내리지 않는다. 동일 RTX 3080/pallet-yolo26 환경에서 N3/PoseFix의 비용 차이는 확인했다. 다만 입력·손실·총학습 및 선택 예산은 완전하게 일치하지 않으므로 정확도·비용의 방법 전체 손익을 조건부로 해석한다.

사람 입력의 어려움 집단에서 위치 오차가 악화하고 큰 오류 복구가 드물다는 결과는 방법의 작동 범위를 제한한다. 출력 박스와 점수 보존은 검출 경로에 개입하지 않았다는 뜻이며, 원래 정확한 코너나 자세를 항상 보존했다는 뜻은 아니다. 코너 개선이 자세 개선으로 전달되는지 직접 확인한 이유도 여기에 있다.

\subsection{자료와 측정의 제한}
직사각형 319장은 재사용 개발자료이고 정사각형은 한 세션·한 치수이다. 직사각형·정사각형의 사람 입력 등급과 참조 코너 가시성은 확보했으나, 등급 판정 기준의 사람 확인은 아직 없고 예측 노출·표본 모집 이력이 독립적으로 통제된 것은 아니다. 두 128장 집합의 frame ID는 전수 일치했으나, 실제 개체 수·선정 이력 및 과거 정사각형 자료와의 전체 노출 관계는 추가 확인 대상이다. 등급 입력 완료를 외부 가림 기준의 승인으로 바꾸지 않았으며, 과거 정사각형150장의7세션 자료는 현119장 평가에 합산하지 않았다.

자세 참조는 수동 코너와 물체 기하의 재구성에 의존한다. 개별 코너 참조의 생성 경로와 작은 차이에 대한 참조 불확실성이 충분히 정량화되지 않았으므로 독립 물리 6D 실측 정확도를 주장하지 않는다. 세션 재표집 구간은 관측 결과의 변동을 설명하지만 독립 세션 검증이나 개발 중 선택 편향 보정을 대신하지 않는다. 기록 영상 사례와 정사각형의 독립 자세 정확도는 실제 참조가 검증된 뒤 별도로 완성해야 한다.

\subsection{작은 개선과 독립 측정의 다음 단계}
치수 입력과 대칭 감독의 약한 추가 이득, 상한 도달 2.62\%, 2D 개선 뒤 이동 악화는 서로 다른 관측이다. 이들만으로 모델 용량 부족이나 상관된 코너 오류가 지배 원인이라고 확정하지 않는다. 같은 동결 특징에서 자세 결합 후보를 선택하는 제한 실험은 탐색적 보충 결과로 분리하며 기존 N3의 방법·결과를 교체하지 않는다. 공동 후보 oracle은 그 bank 안의 선택 여지만 설명하고, 후보별 생성 자세 대신 실제 prediction-only W/D+PnP 최종 출력을 채점해야 한다. 사전에 제한한 GEO/PERM 각 세 seed의 6회 학습을 완료한 보충 탐색에서 GEO의 실사 정규화 코너 오차는 N3/PoseFix보다 악화되었다. RAW/PERM 대비 음의 오차 구간도 좋은 코너 손상 때문에 고정 보존 조건을 충족하지 못했고, 같은 전체 경계의 seed1 비용은 GEO 56.541ms 대 N3 14.433ms였다. 이 고정 후보·특징·목표·예산에서는 새 배포 이득을 지지하지 않으며, 실패를 용량이나 상관 오류의 단일 원인 증명으로 바꾸지 않는다.

실제 동일 상대 자세에서 clean/실물 가림을 취득하고 독립 참조를 설치하는 패킷을 별도로 완성하였다. 현재 8,910장 리프터와 119장 정사각형의 총 9,029개 관측 중 연결된 독립 이동·회전 참조는 0개이므로 독립 pair 평가를 완료했다고 표시하지 않는다. 32개 CAN 중립 명령 구간·8,263프레임은 실제 상대 정지 증거가 아니며, 기존 12장/96점과 후속 72점 저장본도 물리 참조로 승격하지 않는다. 작은 수치 차이를 운용상의 의미 있는 개선으로 판단하려면 참조·교정 오차와 반복성을 함께 계측해야 한다.

